\subsection{可逆子模}

*（我们保持第 5 节的记号。）*

定义 4. 若存在 B 的 A-子模 N 使 M.N = A，则称 B 的 A-子模 M 为可逆的。

例. 若 \(b\) 是 \(\mathbf{B}\) 的可逆元，则 A-模 \(Ab\) 是可逆的，取 \(\mathbf{N} = \mathbf{Ab}^{-1}\) 即可看出。

命题 10. 设 M 是 B 的一个可逆 A-子模。则：

\begin{enumerate}
\def\labelenumi{(\roman{enumi})}
\item
  存在 \(s \in S\) 使 \(As \subset M \subset As^{-1}\) （特别地，\(M\) 是非退化的）。
\item
  A: M 是唯一使 M.N = A 的 B 的 A-子模 N。
\end{enumerate}

若 M.N = A，则 B.M = B.(B.M) ⊃ B.(M.N) = B.A = B，故 M 是非退化的。类似地 N 是非退化的。若 \(t \in S \cap M\) 且 \(u \in S \cap N\) （第 5 节，命题 8），则元素 s = tu 属于 \(S \cap M \cap N\) ，由此 \(Ms \subset M.N = A\) ，因而 \(As \subset M \subset As^{-1}\) 。

另一方面，显然 \(N \subset A: M\) ，由此

\[
\mathbf {A} = \mathbf {M}. \mathbf {N} \subset \mathbf {M}. (\mathbf {A}: \mathbf {M}) \subset \mathbf {A}
\]

以及 M. (A: M) = A；两边乘以 N，推出 A: M = N，证毕。

定理 4. 设 M 是 B 的一个非退化 A-子模。下列性质等价：

\begin{enumerate}
\def\labelenumi{(\alph{enumi})}
\item
  M 是可逆的。
\item
  M 是投射的。
\item
  M 是秩为 1 的投射模。
\item
  M 是有限生成的 A-模，且对 A 的每个极大理想 m，\(A_{m}\) -模 \(M_{m}\) 是单生成的。
\end{enumerate}

首先证明性质 (a)、(b) 与 (c) 的等价性。若 (a) 成立且 N 是使 M.N = A 的 B 的 A-子模，则有关系式

\[
\sum_ {i = 1} ^ {p} m _ {i} n _ {i} = 1 \quad (m _ {i} \in \mathbf {M}, n _ {i} \in \mathbf {N} \text {   for   all   } i).\tag{9}
\]

\[
\boldsymbol {x} \in \mathbf {M},
\]

\[
v _ {i} (x) = n _ {i} x;
\]

\[
\nu_ {i}
\]

\(x = \sum_{i=1}^{n} m_i v_i(x)\) 对所有 \(x \in \mathbf{M}\) 成立；这就证明（《代数》第 II 章 § 2 第 6 节命题 12）\(\mathbf{M}\) 是投射的且由 \(m_i\) 生成；从而 \(\mathbf{M}\) 是有限生成的投射模。

设 \(\mathfrak{m}\) 是 \(\mathbf{A}\) 的一个极大理想；我们证明整数 \(r = \operatorname{rg}_{\mathfrak{m}}(\mathbf{M})\) 等于 1。设 \(S'\) 是 \(S\) 在 \(A_{\mathfrak{m}}\) 中的像；由于 \(S\) 的元素在 \(A\) 中不是零因子，\(S'\) 的元素在 \(A_{\mathfrak{m}}\) 中也不是零因子，因为 \(A_{\mathfrak{m}}\) 是平坦 \(A\) -模（§2 第 4 节定理 1 及第 I 章 §2 第 4 节命题 3）；于是 \(S'^{-1}A_{\mathfrak{m}} \neq 0\) ，并且由于 \(M_{\mathfrak{m}}\) 是自由 \(A_{\mathfrak{m}}\) -模且秩为 \(r\) ，\(S'^{-1}M_{\mathfrak{m}}\) 是自由 \(S'^{-1}A_{\mathfrak{m}}\) -模且秩为 \(r\) 。但若 \(T'\) 是 \(A - \mathfrak{m}\) 在 \(S^{-1}A\) 中的像，则 \(S'^{-1}A_{\mathfrak{m}}\) （相应地，\(S'^{-1}M_{\mathfrak{m}}\) ）典范地等同于 \(T'^{-1}(S^{-1}A)\) （相应地，\(T'^{-1}(S^{-1}M)\) ）（§2 第 3 节命题 7）。而 \(S^{-1}M = B\) （命题 8 (c)），因此 \(T'^{-1}(S^{-1}M)\) 是自由 \(A\) -模，秩为 1，且在 \(T'^{-1}(S^{-1}A)\) 上，这证明了 \(r = 1\) 并给出蕴含关系 (a) ⇒ (c)。

蕴含关系 (c) \(\Rightarrow\) (b) 是平凡的。我们来证明 (b) \(\Rightarrow\) (a)。由假设，存在 M 上的线性型族（不必有限）\((f_{\lambda})_{\lambda \in L}\) 与 M 的元素族 \((m_{\lambda})_{\lambda \in L}\) ，使得对所有 \(x \in M\) ，族 \((f_{\lambda}(x))\) 有有限支撑，且 \(x = \sum_{\lambda \in L} m_{\lambda} f_{\lambda}(x)\) （《代数》第 II 章 §2 第 6 节命题 12）。由于 M 是非退化的，依据第 5 节命题 9，有 \(f_{\lambda}(x) = n_{\lambda} x\) ，其中某个 \(n_{\lambda} \in A: M\) 。取 \(x\) 为 M \(\cap S\) 的一个元素（第 5 节命题 8），可以看出除有限多个指标外必有 \(n_{\lambda} = 0\) ，且 \(\sum_{\lambda \in L} m_{\lambda} n_{\lambda} = 1\) 。这显然蕴含 M. (A: M) = A，即得 (a)。

依据第 3 节的定义 2，(c) 蕴含 (d)。我们来证明其逆。由于 M 是非退化的，它的零化子为零（命题 8 (b)），于是 \(M_{m}\) 的零化子也为零（§ 2 第 4 节公式 (9)）。由于 \(M_{m}\) 假设是单生成的 \(A_{m}\) -模，故它是秩为 1 的自由模，再由第 3 节定理 2 即得 M 是秩为 1 的投射模。

推论. B 的每个可逆 A-子模都是平坦的且是有限表示的。这由定理 4 (c) 得出。

命题 11. 设 M、N 是 B 的两个 A-子模。假设 M 是可逆的。则：

\begin{enumerate}
\def\labelenumi{(\roman{enumi})}
\tightlist
\item
  典范同态 \(\mathbf{M} \otimes_{\mathbf{A}} \mathbf{N} \to \mathbf{M}. \mathbf{N}\) 是双射的。
\end{enumerate}

\[
\text {(ii)} \mathbf {N}: \mathbf {M} = \mathbf {N}. (\mathbf {A}: \mathbf {M}) a n d \mathbf {N} = (\mathbf {N}: \mathbf {M}). \mathbf {M}.
\]

设 \(j\) 是典范单射 \(N \to B\) 。由于 \(M\) 是平坦 A-模（定理 4 的推论），\(1 \otimes j \colon M \otimes_{A} N \to M \otimes_{A} B\) 是单射的。但由 \(B = S^{-1}A\) ，B-模 \(M \otimes_{A} B\) 等于 \(S^{-1}M\) ，因而它等同于 \(B\) ，因为 \(M\) 是非退化的（第 5 节命题 8）。在这一等同之下，\(1 \otimes j\) 的像是 \(M \cdot N\) ，即得 (i)。

令 \(\mathbf{M}' = \mathbf{A} : \mathbf{M}\) 。则显然 \(\mathbf{M}' . \mathbf{N} \subset \mathbf{N} : \mathbf{M}\) 且 \(\mathbf{M} . (\mathbf{N} : \mathbf{M}) \subset \mathbf{N}\) 。另一方面，由 \(\mathbf{M} . \mathbf{M}' = \mathbf{A}\) （命题 10），

\[
\mathbf {N}: \mathbf {M} = \mathbf {M} ^ {\prime}. \mathbf {M}. (\mathbf {N}: \mathbf {M}) \subset \mathbf {M} ^ {\prime}. \mathbf {N}
\]

以及 \(\mathbf{N} = \mathbf{M}.\mathbf{M}'.\mathbf{N}\subset \mathbf{M}.(\mathbf{N}:\mathbf{M})\) ，即得 (ii)。

注. 命题 11 中 (i) 的证明只用到 \(\mathbf{M}\) 是平坦的且非退化的这一事实。

